% Compact standalone background for the new Lerch zero-geometry sections.

\section{Background: the derivative family and its finite zero bound}
\label{lerchboundary:sec:background}

We recall the exact framework from the zero-geometry chapter of the
pinned manuscript \cite{ProveItSnapshot}. The following statements are established inputs to
this continuation; their short proofs are included to make the new
endpoint results readable independently of the repository.

The generalized Stieltjes constants use the convention
\[
 \zeta(s,a)=\frac1{s-1}
 +\sum_{n=0}^{\infty}\frac{(-1)^n}{n!}\gamma_n(a)(s-1)^n,
 \qquad a>0.
\]
For $0\leq\rho<1$, define
\[
 \ell_n(\rho,a)=\sum_{m=0}^{\infty}
      \frac{\rho^m\log^n(a+m)}{a+m},
 \qquad
 \mathcal D_{n,k}^{\rho}(a)=\partial_a^k\ell_n(\rho,a),
 \quad k\geq1.
\]
At $\rho=1$ set $\mathcal D_{n,k}^{1}(a)=\gamma_n^{(k)}(a)$,
where this derivative is with respect to $a$. Only the differentiated
family extends this way: the series defining $\ell_n$ does not
generally converge at $\rho=1$.

Introduce the reciprocal-gamma Appell polynomials
\begin{equation}\label{lerchboundary:eq:Q-background}
 \frac{e^{xt}}{\Gamma(1+t)}
 =\sum_{n=0}^{\infty}Q_n(x)\frac{t^n}{n!},
 \qquad Q_n'(x)=nQ_{n-1}(x).
\end{equation}
The first three are $Q_0=1$, $Q_1=x+\gamma$, and
$Q_2=(x+\gamma)^2-\zeta(2)$.

\begin{proposition}[Appell--Laplace representation]
\label{lerchboundary:prop:laplace-background}
For $n\geq0$, $k\geq1$, $a>0$, and $0\leq\rho\leq1$,
\begin{equation}\label{lerchboundary:eq:laplace-background}
 F_{n,k}^{\rho}(a):=(-1)^{n+k}\mathcal D_{n,k}^{\rho}(a)
 =\int_0^\infty
   \frac{x^k e^{-ax}}{1-\rho e^{-x}}Q_n(\log x)\,dx.
\end{equation}
The integral, and every fixed number of $a$-derivatives, converge
locally uniformly for positive $a$, uniformly in $\rho\in[0,1]$.
In particular $F_{n,k+1}^{\rho}=-\partial_aF_{n,k}^{\rho}$.
\end{proposition}

\begin{proof}
For the Lerch series
$\Phi(\rho,s,a)=\sum_{m\geq0}\rho^m(a+m)^{-s}$,
parameter differentiation gives
$\partial_a^k\Phi(\rho,s,a)=(-1)^k(s)_k\Phi(\rho,s+k,a)$.
At $\rho=1$ the analogous Hurwitz identity follows first from the
convergent series and then by continuation; the Laurent pole at $s=1$
is independent of $a$. Put $s=1+t$ and use the Mellin integral to get
\[
 k!\prod_{j=1}^k(1+t/j)\Phi(\rho,k+1+t,a)
 =\frac1{\Gamma(1+t)}\int_0^\infty
        \frac{x^{k+t}e^{-ax}}{1-\rho e^{-x}}\,dx.
\]
Here $\Phi(1,s,a)$ means $\zeta(s,a)$ in its convergent half-plane.
Extracting the coefficient of $t^n$ proves the formula. Near zero,
$(1-\rho e^{-x})^{-1}\leq(1-e^{-x})^{-1}=O(x^{-1})$,
so the integrand is bounded by a constant times
$x^{k-1}(1+|\log x|^n)$. At infinity, $a$ bounded away from zero
gives exponential decay. These estimates justify the extraction,
all stated $a$-derivatives, and continuity at $\rho=1$.
\end{proof}

\begin{lemma}[Strict real-rootedness of the Appell polynomials]
\label{lerchboundary:lem:Q-roots-background}
For every $n\geq1$, $Q_n$ has $n$ distinct real roots.
\end{lemma}

\begin{proof}
The Weierstrass product \cite{DLMFGamma}
\[
 \frac1{\Gamma(1+t)}
 =e^{\gamma t}\prod_{m=1}^{\infty}(1+t/m)e^{-t/m}
\]
converges locally uniformly. At any finite stage, coefficient
extraction in \eqref{lerchboundary:eq:Q-background} amounts to real
translations and operators $1+\partial_x/m$ applied to $x^n$.
These operators preserve real-rootedness: away from roots of $p$,
new zeros of $p+cp'$ solve $p'/p=-1/c$, whose left side is strictly
decreasing between consecutive poles. Every new zero is simple;
an inherited multiple zero loses one in multiplicity.
Coefficient convergence and closure of monic real-rooted polynomials
give real-rootedness in the infinite-product limit.

To prove simplicity, first remove the first $n-1$ linear factors and
their exponential translation factors from the product. The Appell
polynomial for the remaining tail is still real-rooted by the same
argument. Restore the removed factors. Translations commute with
the differential operators and preserve multiplicity, while the
$n-1$ operators each lower inherited multiplicities and create only
simple new zeros. Thus every final root is simple.
\end{proof}

\begin{proposition}[Global upper bound and persistence]
\label{lerchboundary:prop:bound-background}
For every $n\geq0$, $k\geq1$, and $\rho\in[0,1]$,
$\mathcal D_{n,k}^{\rho}$ has at most $n$ positive zeros counted
with multiplicity. Once it has $n$ distinct positive zeros at an
order $k$, it has exactly $n$ simple positive zeros at every higher
order. With the zeros numbered increasingly,
\[
 a_{j,k}<a_{j,k+1}<a_{j+1,k}\quad(1\leq j<n),
 \qquad a_{n,k}<a_{n,k+1}.
\]
\end{proposition}

\begin{proof}
We use the elementary variation bound for Laplace transforms.
If an integrable signed kernel has $N$ sign changes, multiplication
by $x_0-x$ at a sign-change location removes one sign change.
On its transform $G(a)$ this operation acts as
\[
 (\partial_a+x_0)G(a)
 =e^{-x_0a}\frac{d}{da}\bigl(e^{x_0a}G(a)\bigr).
\]
For $N=0$ the transform has fixed strict sign. Induction, together
with Rolle's theorem counted with multiplicity, proves the bound
of $N$ zeros for general $N$; it also proves nonvanishing of the
transform. The same induction applies to every finite collection
of zeros, so an infinite zero set cannot evade the bound.
In \eqref{lerchboundary:eq:laplace-background}, the positive
weight leaves exactly the $n$ sign changes of $Q_n(\log x)$.
The domination already proved justifies the polynomial moments
used in the variation argument.

Suppose now that $F_{n,k}^{\rho}$ has $n$ distinct zeros. They are
simple by the upper bound. Rolle's theorem gives a derivative zero
between successive zeros. There is also a derivative zero after
the last one: the function has a constant nonzero sign there and
tends to zero at infinity, so it cannot depart from zero and then
return toward zero without a stationary point. The identity
$F_{n,k+1}^{\rho}=-\partial_aF_{n,k}^{\rho}$ supplies at least $n$
distinct zeros at the next order. The upper bound supplies exact
completeness and simplicity, and excludes any extra zero or
coincidence with an old simple zero. Iterate this argument.
The limit at infinity follows by dominated convergence in the
Laplace representation, using any fixed positive reference $a$.
\end{proof}

These are the manuscript's existing positivity and variation arguments.
The new theorems below use their global multiplicity bound to turn
local perturbation information and finitely certified signs into
complete zero classifications.
